\documentclass[12pt,a4paper]{article}
\usepackage[utf8]{inputenc}
\usepackage[T1]{fontenc}
\usepackage[french]{babel}
\usepackage{geometry}
\usepackage{titlesec}
\usepackage{booktabs}
\usepackage{array}
\usepackage{tabularx}
\usepackage{enumitem}
\usepackage{graphicx}
\usepackage{hyperref}
\usepackage{fancyhdr}
\usepackage{xcolor}
\usepackage{parskip}

\geometry{top=2.5cm, bottom=3cm, left=2.5cm, right=2.5cm}

\definecolor{primary}{RGB}{0,70,127}
\hypersetup{
  colorlinks=true,
  linkcolor=primary,
  urlcolor=primary
}

\pagestyle{fancy}
\fancyhf{}
\renewcommand{\headrulewidth}{0.4pt}
\renewcommand{\footrulewidth}{0.4pt}
\fancyhead[L]{\small UE : ANI-IA 4067 (AIA4)}
\fancyhead[R]{\small ENSPY -- Génie Informatique}
\fancyfoot[L]{\small NEUSSI NJIETCHEU Patrice Eugene -- 24P820}
\fancyfoot[C]{\thepage}
\fancyfoot[R]{\includegraphics[height=0.8cm]{enspy.png}}

\fancypagestyle{plain}{
    \fancyhf{}
    \fancyfoot[L]{\small NEUSSI NJIETCHEU Patrice Eugene -- 24P820}
    \fancyfoot[C]{\thepage}
    \fancyfoot[R]{\includegraphics[height=0.8cm]{enspy.png}}
    \renewcommand{\footrulewidth}{0.4pt}
    \renewcommand{\headrulewidth}{0pt}
}

\titleformat{\section}{\large\bfseries\color{primary}}{}{0em}{\thesection\quad}[\titlerule]
\titleformat{\subsection}{\normalsize\bfseries}{}{0em}{\thesubsection\quad}

\setlist[itemize]{noitemsep, topsep=4pt, leftmargin=1.5em}
\setlist[enumerate]{noitemsep, topsep=4pt, leftmargin=1.8em}

\begin{document}

%---------------------------------------------------------------
%  PAGE DE TITRE
%---------------------------------------------------------------
\begin{titlepage}
  \centering
  \includegraphics[width=4cm]{enspy.png}\\[0.5cm]
  {\Large\bfseries ÉCOLE NATIONALE SUPÉRIEURE POLYTECHNIQUE DE YAOUNDÉ\par}
  \vspace{0.2cm}
  {\normalsize Département de Génie Informatique\par}
  {\normalsize Cycle des Ingénieurs du Numérique -- Filière : AIA4\par}
  \vspace{2cm}
  
  {\LARGE\bfseries\color{primary} RÉCAPITULATIF DES PROJETS PRATIQUES\par}
  \vspace{0.5cm}
  {\large Bilans de Progression, Objectifs Atteints et Difficultés Rencontrées\par}
  \vspace{0.2cm}
  {\large\bfseries UE : ANI-IA 4067 INTELLIGENCE ARTIFICIELLE (AIA4)\par}
  \vspace{2.5cm}
  
  \begin{tabular}{rl}
    \textbf{Rédigé par :} & \textbf{NEUSSI NJIETCHEU Patrice Eugene} \\[4pt]
    \textbf{Matricule :}    & \textbf{24P820} \\[4pt]
    \textbf{Classe :}       & \textbf{AIA4} (Intelligence Artificielle) \\[4pt]
    \textbf{Sous la direction de :} & \textbf{M. Bitha Junior} (Enseignant ENSPY) \\
  \end{tabular}
  \vfill
  {\small Bilan d'avancement au 21 mai 2026\par}
\end{titlepage}

%---------------------------------------------------------------
%  TABLE DES MATIÈRES
%---------------------------------------------------------------
\tableofcontents
\newpage

%---------------------------------------------------------------
%  INTRODUCTION
%---------------------------------------------------------------
\section{Bilan de Progression et de Réalisation}

Ce rapport dresse le bilan d'avancement des trois projets d'Intelligence Artificielle qui nous ont été alloués dans le cadre de l'UE ANI-IA 4067 à la date du **21 mai 2026**. Pour chaque projet, nous détaillons les objectifs déjà atteints, les difficultés d'implémentation et algorithmiques actuellement rencontrées, ainsi que le pourcentage de réalisation estimé.

\bigskip

\begin{table}[H]
    \centering
    \caption{Tableau de synthèse de progression des projets au 21 mai}
    \vspace{0.2cm}
    \begin{tabularx}{\linewidth}{lXcc}
        \toprule
        \textbf{N°} & \textbf{Intitulé du Projet} & \textbf{État de Réalisation} & \textbf{Pourcentage} \\
        \midrule
        1 & Recommandation par GNN (LightGCN) & En cours d'intégration web & 80\% \\
        2 & Moteur de Recherche Visuelle (CBIR) & En cours d'optimisation serverless & 85\% \\
        3 & Classification Comportementale (Banque) & En cours d'assemblage d'ensemble & 75\% \\
        \bottomrule
    \end{tabularx}
\end{table}

\newpage

%---------------------------------------------------------------
%  PROJET 1
%---------------------------------------------------------------
\section{Projet 1 -- Système de Recommandations Graph Neural Networks (LightGCN)}

\subsection{Objectifs Atteints}
\begin{itemize}
    \item \textbf{Préparation du Graphe Bipartite (100\%) :} Nettoyage du corpus d'interaction MovieLens-100k, structuration sous forme de matrice d'adjacence globale et calcul de la normalisation symétrique ($\tilde{A} = D^{-1/2} A D^{-1/2}$).
    \item \textbf{Modélisation LightGCN (100\%) :} Implémentation du modèle de propagation à 3 couches sous PyTorch et codage de la boucle d'apprentissage par perte de classement pairwise (BPR Loss).
    \item \textbf{Entraînement et Validation (100\%) :} Validation hors-ligne du modèle démontrant un Recall@10 de 0.0825 et un NDCG@10 de 0.1421.
\end{itemize}

\subsection{Difficultés Rencontrées et Actions en Cours}
\begin{itemize}
    \item \textbf{Sparsité extrême de la matrice d'interactions :} La rareté des notations (supérieure à 93\%) rend les profils d'utilisateurs "froids" très bruités.
    \textit{Action :} Le passage de messages d'ordre supérieur de LightGCN permet de lisser les représentations, mais nous étudions l'intégration d'un profil synthétique de départ pour surmonter le problème de cold-start.
    \item \textbf{Architecture de routage Django-Vis.js :} La modélisation sous forme de réseau dynamique Vis.js dans l'interface web nécessite le transfert efficace des liaisons du graphe.
    \textit{Action :} Nous terminons l'écriture du script de sérialisation JSON pour envoyer les nœuds et les arêtes d'interaction au template HTML de Django.
\end{itemize}

\subsection{Pourcentage de Réalisation}
\textbf{80\%} -- Le cœur de la modélisation mathématique et du modèle IA est validé ; l'intégration web et le déploiement serverless sont en cours de finalisation.

\newpage

%---------------------------------------------------------------
%  PROJET 2
%---------------------------------------------------------------
\section{Projet 2 -- Moteur de Recherche Visuel}

\subsection{Objectifs Atteints}
\begin{itemize}
    \item \textbf{Indexation du catalogue (100\%) :} Constitution et redimensionnement d'une base statique de 150 vêtements, et écriture du pipeline d'indexation hors-ligne.
    \item \textbf{Extraction de Features Sémantiques (100\%) :} Chargement du CNN MobileNetV2 et extraction de descripteurs profonds de dimension 1280.
    \item \textbf{Descripteurs Spatiaux et Colorimétriques (100\%) :} Code d'extraction d'histogramme joint 3D Hue-Saturation-Value (128 dimensions) et grille de nuances de gris normalisée Z-score (256 dimensions) terminé.
\end{itemize}

\subsection{Difficultés Rencontrées et Actions en Cours}
\begin{itemize}
    \item \textbf{Limite de build sur l'architecture sans état (Vercel) :} L'importation de PyTorch et le téléchargement des poids de MobileNetV2 dépassent la taille limite de 15MB compressés allouée par Vercel aux fonctions serverless.
    \textit{Action :} Nous concevons un algorithme de repli (fallback) hybride. Si PyTorch ne peut pas être chargé sur le serveur, l'application basculera automatiquement sur l'extraction locale HSV + Spatial, ce qui permet de bypasser l'installation de PyTorch et de conserver un temps de recherche minimal sous les 100ms.
\end{itemize}

\subsection{Pourcentage de Réalisation}
\textbf{85\%} -- Les deux moteurs de recherche (sémantique et colorimétrique) sont fonctionnels en local. Le script de repli dynamique est en cours d'intégration dans l'interface de recherche Django.

\newpage

%---------------------------------------------------------------
%  PROJET 3
%---------------------------------------------------------------
\section{Projet 3 -- Prédiction du Comportement Client (Banque)}

\subsection{Objectifs Atteints}
\begin{itemize}
    \item \textbf{Analyse Exploratoire \& Prétraitement (100\%) :} EDA complète du dataset UCI Bank Marketing, encodage One-Hot des données catégorielles et standardisation robuste des variables financières.
    \item \textbf{Modélisation Individuelle (100\%) :} Entraînement de la Régression Logistique et de la Forêt Aléatoire avec ajustement des poids de classe.
    \item \textbf{Réseau de Neurones MLP (100\%) :} Codage et entraînement sous PyTorch du perceptron multicouche à 3 couches linéaires avec Dropout.
\end{itemize}

\subsection{Difficultés Rencontrées et Actions en Cours}
\begin{itemize}
    \item \textbf{Déséquilibre prononcé des cibles (11.7\% d'adhésion) :} Les modèles tendent à maximiser l'accuracy au détriment du rappel de la classe d'intérêt.
    \textit{Action :} Nous mettons au point une architecture d'ensemble par vote pondéré (Soft Voting) regroupant la régression logistique, la forêt aléatoire et le MLP pour lisser les scores de probabilité et cibler le compromis optimal.
    \item \textbf{Encapsulation PyTorch dans scikit-learn :} Le framework d'ensemble de scikit-learn nécessite que le modèle PyTorch implémente les APIs `fit` et `predict\_proba`.
    \textit{Action :} Nous développons une classe enveloppe personnalisée `PyTorchMLPClassifier` qui hérite des classes de base de scikit-learn pour s'insérer proprement dans le `VotingClassifier`.
\end{itemize}

\subsection{Pourcentage de Réalisation}
\textbf{75\%} -- L'analyse statistique et l'entraînement des modèles de base sont faits. Nous finalisons l'intégration de la classe enveloppe d'ensemble et l'interface de prédiction en lot (batch) Django.

\end{document}
